% from ega0-1-fr.tex
\oldpage[0\textsubscript{I}]{21}

\section{Espaces irréductibles. Espaces noethériens}
\label{section:0.2-fr}

\subsection{Espaces irréductibles}
\label{subsection:0.2.1-fr}

\begin{env}[2.1.1]
\label{0.2.1.1-fr}
On dit qu'un espace topologique $X$ est \emph{irréductible} s'il est non vide
et s'il n'est pas réunion de deux sous-espaces fermés distincts de $X$. Il
revient au même de dire que $X\neq\varnothing$ et que l'intersection de deux
ouverts (et par suite d'un nombre fini d'ouverts) non vides de $X$ est non
vide, ou que tout ouvert non vide est partout dense, ou que toute partie
fermée $\neq X$ est \emph{rare}, ou enfin que tout ouvert de $X$ est
\emph{connexe}.
\end{env}

\begin{env}[2.1.2]
\label{0.2.1.2-fr}
Pour qu'un sous-espace $Y$ d'un espace topologique $X$ soit irréductible, il
faut et il suffit que son adhérence $\overline Y$ soit irréductible. En
particulier, tout sous-espace qui est l'adhérence $\overline{\{x\}}$ d'un
sous-espace réduit à un point est irréductible; nous exprimerons la relation
$y\in\overline{\{x\}}$ (équivalente à
$\overline{\{y\}}\subset\overline{\{x\}}$) en disant que $y$ est
\emph{spécialisation de $x$} ou que $x$ est une \emph{générisation de $y$}.
Lorsqu'il existe dans un espace irréductible $X$ un point $x$ tel que
$X=\overline{\{x\}}$, nous dirons que $x$ est \emph{point générique} de $X$.
Tout ouvert non vide de $X$ contient alors $x$ et tout sous-espace contenant
$x$ admet $x$ pour point générique.
\end{env}

\begin{env}[2.1.3]
\label{0.2.1.3-fr}
Rappelons qu'on appelle \emph{espace de Kolmogoroff} un espace topologique $X$
vérifiant l'axiome de séparation~:

\smallskip
\noindent
\textup{(T$_0$)} Si $x\neq y$ sont deux points quelconques de $X$, il existe
un ensemble ouvert contenant l'un des points $x$, $y$ et non l'autre.

Si un espace de Kolmogoroff irréductible admet un point générique, il n'en
admet \emph{qu'un seul} puisqu'un ouvert non vide contient tout point
générique.

Rappelons qu'un espace topologique $X$ est dit \emph{quasi-compact} si, de tout
recouvrement ouvert de $X$, on peut extraire un recouvrement fini de $X$ (ou,
ce qui revient au même, si toute famille filtrante décroissante d'ensembles
fermés non vides a une intersection non vide). Si $X$ est un espace
quasi-compact, toute partie fermée non vide $A$ de $X$ contient un ensemble
fermé non vide \emph{minimal} $M$, car l'ensemble des parties fermées non vides
de $A$ est inductif pour la relation $\supset$; si en outre $X$ est un espace
de Kolmogoroff, $M$ est nécessairement réduite à un seul point (ou, comme on
dit par abus de langage, est un \emph{point fermé}).
\end{env}

\begin{env}[2.1.4]
\label{0.2.1.4-fr}
Dans un espace irréductible $X$, tout sous-espace ouvert non vide $U$ est
irréductible, et si $X$ admet un point générique $x$, $x$ est aussi point
générique de $U$.

Soit $(U_\alpha)$ un recouvrement (dont l'ensemble d'indices est non vide) d'un
espace topologique $X$, formé d'ouverts non vides; pour que $X$ soit
irréductible, il faut et il suffit que $U_\alpha$ soit irréductible pour tout
$\alpha$, et que $U_\alpha\cap U_\beta\neq\varnothing$ quels que soient
$\alpha$, $\beta$. La condition est évidemment nécessaire; pour voir qu'elle
est suffisante, il suffit de prouver que si $V$ est un ouvert non vide de $X$,
$V\cap U_\alpha$ est non vide pour tout $\alpha$, car alors $V\cap U_\alpha$
est dense dans $U_\alpha$ pour tout $\alpha$, et par suite $V$ est dense dans
$X$. Or, il y a au moins un indice $\gamma$ tel que
$V\cap U_\gamma\neq\varnothing$, donc $V\cap U_\gamma$ est dense dans
$U_\gamma$, et comme, pour tout $\alpha$,
$U_\alpha\cap U_\gamma\neq\varnothing$, on a aussi
$V\cap U_\alpha\cap U_\gamma\neq\varnothing$.
\end{env}

